% ============================================================
%  GENERADO POR generar_suplemento.py -- NO EDITAR A MANO
%  Fuente: experimento/resultados/c2_ne.json
%  Corrida: 2026-09-13T21:38:26Z  code_sha256_16=2f4ecfd3aa636719
% ============================================================
\documentclass[11pt]{article}
\usepackage[T1]{fontenc}
\usepackage{lmodern}
\usepackage[a4paper,margin=2.5cm]{geometry}
\usepackage{amsmath,amssymb}
\usepackage{booktabs}
\usepackage[authoryear,round]{natbib}
\renewcommand{\thetable}{S\arabic{table}}
\title{Supplementary material\\[2pt]\large Cohort and
generation-interval measurements}
\author{Ra\'ul-Antonio Ja\'en-Garc\'ia}
\date{}
\begin{document}
\maketitle

\noindent This document accompanies \emph{Extrinsic Is Not One Thing:
Recipient Exchangeability in Evolutionary Systems}. It reports the
cohort and generation-interval measurements referred to in the
substrate-validation section of the main text. Every number below is
generated from the same result file that produces the corresponding
statements in the main text, by the script
\texttt{generar\_suplemento.py} in the repository.

\section*{S1. What was measured, and why}

The effective population size reported in the main text is estimated
from the decay of neutral marker variance. Comparing that estimate
against Wright's $N_e = 4N/(V_k+2)$ requires the variance in offspring
number $V_k$, and comparing it against the overlapping-generation
correction of \citet{hill1972} requires in addition the generation
interval $T$.
Both are properties of cohorts rather than of instantaneous state, so
they are measured by following identified cohorts to extinction.

Individuals are identified by the pair (slot, birth time). The
substrate reuses array slots when an individual dies, so identifying
cohort members by slot alone counts several successive occupants as
one individual; that error produced offspring counts of 45 and 188 per
individual before it was corrected. The cohort loop also runs for the
full horizon at every population size rather than stopping when the
focal cohort is exhausted, so that the comparison across sizes is over
the same number of generations of drift.

\section*{S2. Generation interval and offspring-number variance}

\begin{table}[ht]
\centering\small
\caption{Cohort measurements at three census sizes. $\bar{k}$ is mean
offspring number per individual followed to extinction, $V_k$ its
variance, $T$ the generation interval in substrate steps, and
$T/\mathrm{lifetime}$ the interval expressed in mean lifetimes.}
\label{tab:s_cohortes}
\begin{tabular}{rrrrrrr}
\toprule
$R$ & $N$ & $\bar{k}$ & $V_k$ & $T$ & mean lifetime & $T/\mathrm{lifetime}$\\
\midrule
250 & 368.7 & 1.362 & 1.868 & 109.98 & 80.40 & 1.368\\
500 & 737.2 & 1.380 & 1.965 & 109.90 & 80.28 & 1.369\\
1\,000 & 1\,474.8 & 1.369 & 1.844 & 109.93 & 80.29 & 1.369\\
\bottomrule
\end{tabular}
\end{table}

Two features of Table~\ref{tab:s_cohortes} carry the argument of the
main text. First, the generation interval is \emph{flat}: $T$ varies
between 109.90 and 109.98 steps while $N$ grows from
369 to 1\,475, a spread of
0.07\% across a fourfold change
in census size. Second, $V_k$ is likewise approximately constant, at
1.84 to 1.96, with $\bar{k}$ between
1.362 and 1.380.

\section*{S3. $V_k$ does not explain the sublinear scaling}

\begin{table}[ht]
\centering\small
\caption{Measured $N_e$ against Wright's prediction and against the
overlapping-generation correction. The correction moves the prediction
in the wrong direction at every size.}
\label{tab:s_ne}
\begin{tabular}{rrrrrr}
\toprule
$N$ & $N_e$ measured & $4N/(V_k+2)$ & ratio & Hill $4N_bT/(V_k+2)$ & ratio\\
\midrule
368.7 & 197.8 & 381.3 & 0.519 & 521.7 & 0.379\\
737.2 & 325.5 & 743.8 & 0.438 & 1\,018.3 & 0.320\\
1\,474.8 & 461.8 & 1\,534.8 & 0.301 & 2\,101.3 & 0.220\\
\bottomrule
\end{tabular}
\end{table}

With $V_k$ near two, Wright's expression predicts $N_e \approx N$. The
measured value is 0.30 to 0.52 of that, and the
ratio \emph{falls} as $N$ grows, which is the sublinear scaling
reported in the main text. Since $V_k$ is flat across the same range,
it cannot be what produces the scaling, and we have no mechanism that
does.

The overlapping-generation correction of \citet{hill1972}, $N_e =
4 N_b T / (V_k + 2)$, is the natural candidate, since this substrate has
continuous reproduction and overlapping generations. It does not help.
Because $T$ is constant at about 110 steps rather
than growing with $N$, the correction multiplies the prediction by a
roughly fixed factor, raising it further above a measurement that was
already below it: the ratio moves from
0.30--0.52 to 0.22--0.38,
that is, further from one at every size. We report the hypothesis as
falsified rather than adjusting it.

\section*{S4. The two estimators of $N_e$}

\begin{table}[ht]
\centering\small
\caption{The drift-based and mutation--drift-balance estimators of
$N_e$, and their ratio. Agreement improves as $N$ grows.}
\label{tab:s_estimadores}
\begin{tabular}{rrrr}
\toprule
$N$ & $N_e$ from drift & $N_e$ from mutation--drift & ratio\\
\midrule
368.7 & 197.8 & 91.7 & 2.158\\
737.2 & 325.5 & 160.0 & 2.034\\
1\,474.8 & 461.8 & 374.1 & 1.234\\
\bottomrule
\end{tabular}
\end{table}

The two estimators are independent: one reads the decay of neutral
variance, the other its equilibrium level under recurrent mutation.
They agree to within a factor of 1.2 to 2.2,
and the agreement improves monotonically with $N$. We report this as
partial concordance rather than confirmation: a factor of two between
two estimators of the same quantity is a real discrepancy, not a
rounding difference, and it is consistent with the sublinear scaling
of \S S3 being a property of the substrate that neither estimator
fully captures.

\section*{S5. Provenance}

\begin{tabular}{ll}
\toprule
Result file & \texttt{c2\_ne.json}\\
Run completed & 2026-09-13T21:38:26Z\\
Source hash (\texttt{code\_sha256\_16}) & \texttt{2f4ecfd3aa636719}\\
Python & 3.10.5\\
NumPy & 1.26.4\\
Seeds per point & 16\\
Burn-in steps & 5\,000\\
\bottomrule
\end{tabular}

\bibliographystyle{elsarticle-harv}
\bibliography{references}
\end{document}
